\documentclass[11pt]{article}

\usepackage{amsmath,amssymb,amsthm,mathtools}
\usepackage[T1]{fontenc}
\usepackage{lmodern}
\usepackage[margin=1in]{geometry}
\usepackage{enumitem}
\usepackage{hyperref}
\usepackage{microtype}

\hypersetup{
  colorlinks=true,
  linkcolor=blue,
  citecolor=blue,
  urlcolor=blue
}

\newtheorem{theorem}{Theorem}[section]
\newtheorem{proposition}[theorem]{Proposition}
\newtheorem{corollary}[theorem]{Corollary}
\newtheorem{definition}[theorem]{Definition}
\newtheorem{lemma}[theorem]{Lemma}
\newtheorem{remark}[theorem]{Remark}

\newcommand{\Z}{\mathbb Z}
\newcommand{\Acal}{\mathcal A}
\newcommand{\Word}{\mathcal W}
\newcommand{\Pad}{\mathsf{Pad}}
\newcommand{\tauv}{\tau}
\newcommand{\Ev}{\mathsf E}

\title{A compact reproof of the 077 fixed-delta tail-length reduction}
\author{}
\date{March 2026}

\begin{document}
\maketitle

\begin{abstract}
This note removes the standalone 077 support theorem by proving directly the
exact statement used later in 092.

The reproof is short once the bridge platform is written explicitly in cleaned
form. The bridge proposition already used in 092 supplies two decisive facts:
\begin{enumerate}[label=\textup{(\arabic*)},leftmargin=2.5em]
\item on a full chain with boundary label $\delta$, the current event class is
      an explicit function of $(\beta,\delta)$;
\item whenever a forward successor chain exists, its boundary label is
      $\delta+1$.
\end{enumerate}
So a lift starting at $\delta$ emits a deterministic chain block
$\Word(\delta)$, then $\Word(\delta+1)$, then $\Word(\delta+2)$, and so on,
until terminality or forever.

Therefore the padded future current-event word of an actual lift is determined
by exactly two pieces of data:
\begin{itemize}[leftmargin=2em]
\item its starting label $\delta$,
\item its remaining full-chain tail length.
\end{itemize}
Equivalently, at fixed realized $\delta$, two actual lifts can differ only by
remaining tail length. The old 080 no-mixed-delta reduction then becomes an
immediate corollary.
\end{abstract}

\section{Scope and imported bridge platform}

Fix an odd modulus $m$ in the accepted D5 regime and the best-seed channel
\[
R1 \to H_{L1}.
\]

We use the standard cleaned definitions from 092.

\begin{definition}[full chains, actual lifts, tail length]
A \emph{full chain} is a complete boundary block between consecutive visits to
the $\Theta=2$ section.

An \emph{actual lift} of a realized boundary label $\delta_0$ is an actual
accessible full-chain occurrence with boundary label $\delta_0$.

The \emph{remaining tail length} $\tauv(C)$ of an actual lift $C$ is the number
of full-chain successors available after $C$ before terminality. If successors
continue forever, then $\tauv(C)=\infty$.
\end{definition}

\begin{definition}[endpoint class and padding]
Fix the standard accepted terminal padding convention from 092 for future
current-event words. Two actual lifts lie in the same \emph{endpoint class}
when they determine the same padded future current-event word.

Write $\Pad$ for the common terminal padding tail supplied by that convention.
The corresponding boundary right-congruence class is denoted by $\rho$.
\end{definition}

Let
\[
\Acal=
\{\mathrm{wrap},\mathrm{carry\_jump},\mathrm{other}_{1000},
  \mathrm{other}_{0010},\mathrm{flat}\}
\]
be the current-event alphabet.

The present note imports only the bridge platform already written explicitly in
cleaned form in Proposition~4.1 of 092:
\begin{enumerate}[label=\textup{(\arabic*)},leftmargin=2.5em]
\item every full chain starts at the common boundary clock value
      $\beta=m-2$, and during that chain the clock evolves by
      \[
      \beta \mapsto \beta-1 \pmod m;
      \]
\item there is an explicit current-event law
      \[
      \Ev:\Z/m\Z \times \Z/m^2\Z \longrightarrow \Acal,
      \qquad
      (\beta,\delta) \longmapsto \epsilon_4(\beta,\delta),
      \]
      which is independent of any component tag;
\item if a full chain $C'$ is the successor of a full chain $C$, then
      \[
      \delta(C') = \delta(C)+1 \pmod{m^2}.
      \]
\end{enumerate}

Item \textup{(2)} is exactly the explicit event law carried by the componentwise
bridge platform. Item \textup{(3)} is the componentwise odometer splice law.

\section{Deterministic chain blocks}

\begin{definition}[the one-chain block at label $\delta$]
For $\delta\in\Z/m^2\Z$, define the canonical one-chain event block
\[
\Word(\delta)
=\bigl(\Ev(m-2,\delta),\Ev(m-3,\delta),\dots,\Ev(0,\delta),\Ev(m-1,\delta)\bigr)
\in \Acal^m.
\]
This is the current-event word read while the chain clock runs once through its
common boundary-to-boundary cycle.
\end{definition}

\begin{lemma}[the current chain word is determined by $\delta$]
Let $C$ be any actual lift with boundary label $\delta$. Then the current-event
word emitted along the chain $C$ is exactly $\Word(\delta)$.
\end{lemma}

\begin{proof}
Every full chain starts at the common boundary clock value $\beta=m-2$, and
inside that chain the clock decreases deterministically by $1$ modulo $m$.
The emitted current event at each position is therefore read from the same
explicit law $\Ev(\beta,\delta)$, with the same ordered list of clock values
\[
m-2,m-3,\dots,0,m-1.
\]
So the full chain emits exactly the block $\Word(\delta)$.
\end{proof}

\begin{proposition}[block stack determined by start label and tail length]
Let $C$ be an actual lift with boundary label $\delta$.
\begin{enumerate}[label=\textup{(\arabic*)},leftmargin=2.5em]
\item If $\tauv(C)=t<\infty$, then the forward chain labels are
      \[
      \delta,\ \delta+1,\ \dots,\ \delta+t,
      \]
      where the chain at $\delta+t$ is terminal, and the padded future
      current-event word of $C$ is
      \[
      \Word(\delta)\,\Word(\delta+1)\,\cdots\,\Word(\delta+t)\,\Pad.
      \]
\item If $\tauv(C)=\infty$, then the padded future current-event word of $C$
      is the infinite concatenation
      \[
      \Word(\delta)\,\Word(\delta+1)\,\Word(\delta+2)\,\cdots.
      \]
\end{enumerate}
In particular, the padded future current-event word is completely determined by
$\delta$ together with $\tauv(C)$.
\end{proposition}

\begin{proof}
Start with the current chain. By Lemma~2.2, it emits $\Word(\delta)$.

If no successor exists, then $\tauv(C)=0$, the current chain is terminal, and
by the fixed padding convention the future word is
\[
\Word(\delta)\Pad.
\]
So the claim holds in the base case.

Now suppose a successor chain exists. By the odometer splice law from the
imported bridge platform, its boundary label is $\delta+1$. By Lemma~2.2 again,
that successor emits the block $\Word(\delta+1)$. Repeating the same argument
for each further successor yields the entire stack
\[
\Word(\delta),\Word(\delta+1),\Word(\delta+2),\dots
\]
until the first terminal chain, if one occurs, or forever if no terminal chain
occurs.

Thus a finite tail of length $t$ gives exactly the finite concatenation through
$\delta+t$ followed by the common padding suffix $\Pad$, while an infinite tail
gives the infinite concatenation. No other local future-law choice remains.
\end{proof}

\section{The reproof replacing 077}

\begin{theorem}[fixed-$\delta$ ambiguity reduces to tail length]
Let $C_1$ and $C_2$ be actual lifts of the same realized boundary label
$\delta$.

Then the following hold:
\begin{enumerate}[label=\textup{(\arabic*)},leftmargin=2.5em]
\item if $\tauv(C_1)=\tauv(C_2)$, then $C_1$ and $C_2$ have the same padded
      future current-event word;
\item $C_1$ and $C_2$ have the same padded future current-event word if and
      only if they lie in the same endpoint class;
\item $C_1$ and $C_2$ lie in the same endpoint class if and only if
      $\rho(C_1)=\rho(C_2)$.
\end{enumerate}

Equivalently, at fixed realized $\delta$, the padded future current-event word,
the endpoint class, and the boundary right-congruence class $\rho$ are all
determined by the remaining full-chain tail length. There is no further local
future-law ambiguity at fixed realized $\delta$.
\end{theorem}

\begin{proof}
By Proposition~2.3, the padded future current-event word of an actual lift is a
function only of its starting label $\delta$ and its tail length $\tauv$.
Therefore, for two lifts starting at the same realized label $\delta$, equality
of tail length is sufficient for equality of the padded future word.

Items \textup{(2)} and \textup{(3)} are merely the definitions of endpoint class
and of the right-congruence label $\rho$ attached to the padded future word.
\end{proof}

\begin{corollary}[first divergence occurs at the shorter terminal row]
Let $C_1$ and $C_2$ be actual lifts of the same realized label $\delta$, and
assume $\tauv(C_1)\neq\tauv(C_2)$. Set
\[
t = \min\{\tauv(C_1),\tauv(C_2)\} < \infty.
\]
Then the later realized label
\[
\delta+t
\]
is mixed-status: one actual lift at that label is terminal and another actual
lift at that same label continues.
\end{corollary}

\begin{proof}
By Proposition~2.3, both lifts follow the same deterministic block stack
through the chain labels
\[
\delta,\ \delta+1,\ \dots,\ \delta+t.
\]
At the last of these labels, namely $\delta+t$, the lift with shorter tail has
no further successor, while the other lift still has at least one more
successor. So $\delta+t$ is realized with both a terminal lift and a continuing
lift.
\end{proof}

\begin{corollary}[080 becomes immediate]
If no realized boundary label is mixed-status, then fixed realized $\delta$
determines tail length.

Consequently, fixed realized $\delta$ determines endpoint class and the padded
future current-event word.
\end{corollary}

\begin{proof}
If fixed realized $\delta$ did \emph{not} determine tail length, then two lifts
of that $\delta$ would have different tail lengths, and Corollary~3.2 would
produce a later mixed-status realized label. So no mixed-status implies
realization-independent tail length at fixed $\delta$.

The endpoint-class and padded-word conclusions then follow from Theorem~3.1.
\end{proof}

\begin{corollary}[how 097 is used in the final globalization proof]
Assume every actual lift has infinite tail length. Then fixed realized
$\delta$ determines:
\begin{enumerate}[label=\textup{(\arabic*)},leftmargin=2.5em]
\item the padded future current-event word,
\item the endpoint class,
\item the boundary right-congruence class $\rho$.
\end{enumerate}
Hence $\rho=\rho(\delta)$ globally.
\end{corollary}

\begin{proof}
Under the hypothesis, any two lifts of the same realized $\delta$ have the same
tail length, namely $\infty$. Apply Theorem~3.1.
\end{proof}

\section{What 097 changes in the cleaned package}

\begin{corollary}[accepted-support list after 097]
Relative to the cleaned 092 package and the already completed reproof notes
095--096, the old 077 theorem no longer needs to remain in accepted-support
form.

So after 097, the cleaned odd-$m$ D5 package no longer imports 077, 079, or
081 as separate support theorems. The remaining imported inputs are the
componentwise bridge platform and the localized structural first-exit block
used inside Sections~3 and~4 of 092.
\end{corollary}

\begin{remark}[what this note does and does not use]
This note is a pure theorem-level reproof. Unlike the older 077 support note,
it does \emph{not} rely on frozen actual-union checks for small moduli. The
proof needs only the explicit bridge law written in cleaned form and the
standard 092 definitions.

In particular, the old 080 reduction is now compressed to Corollary~3.3 inside
this same note.
\end{remark}

\end{document}
